\documentclass[11pt,a4paper]{article}
\usepackage[margin=25mm]{geometry}
\usepackage[T1]{fontenc}
\usepackage{lmodern}
\usepackage{amsmath,amssymb,amsthm}
\usepackage{microtype}
\usepackage{xcolor}
\usepackage{fancyhdr}
\usepackage[colorlinks=true,linkcolor=blue!45!black]{hyperref}
\hypersetup{pdftitle={Conditional finiteness for x cubed y squared = z to the fourth + 1},pdfsubject={A self-contained proof assuming the abc conjecture}}
\DeclareMathOperator{\rad}{rad}
\newtheorem{theorem}{Theorem}
\newtheorem*{conjecture}{The abc conjecture}
\pagestyle{fancy}
\fancyhf{}
\fancyhead[L]{\small Conditional Diophantine finiteness}
\fancyhead[R]{\small $x^3y^2=z^4+1$}
\fancyfoot[C]{\small\thepage}
\setlength{\headheight}{14pt}
\setlength{\parindent}{0pt}
\setlength{\parskip}{5pt}
\renewcommand{\headrulewidth}{0.3pt}
\begin{document}
\begin{center}
{\LARGE\bfseries Conditional finiteness for}\\[6pt]
{\Large $\boldsymbol{x^3y^2=z^4+1}$}\\[7pt]
{\small A self-contained proof assuming the $abc$ conjecture}
\end{center}

The equation has only finitely many integer solutions \emph{if the
$abc$ conjecture holds}. All steps below are elementary once that conjecture
is assumed. This proves the implication; it does not establish finiteness
independently of the conjecture.

\section*{1. Precise statement and proof}
For a positive integer $n$, its \emph{radical} $\rad(n)$ is the product of
its distinct prime divisors, with $\rad(1)=1$. In particular,
$\rad(n)\le n$ and $\rad(mn)\le\rad(m)\rad(n)$ for positive integers $m,n$.
These inequalities follow directly by counting each prime at most once.

\begin{conjecture}
For every $\varepsilon>0$ there is a constant $K_\varepsilon\ge1$ such that
all pairwise coprime positive integers $a,b,c$ satisfying $a+b=c$ obey
\[
c\le K_\varepsilon\,\rad(abc)^{1+\varepsilon}.
\]
\end{conjecture}

\begin{theorem}
Assume the $abc$ conjecture, and let $K=K_{1/6}$. Every integer solution of
$x^3y^2=z^4+1$ satisfies
\[
1\le x\le K^{8/3},\qquad 1\le |y|\le K^4,\qquad |z|\le K^2.
\]
Consequently the equation has only finitely many integer solutions.
\end{theorem}

\begin{proof}
The right-hand side is positive, so $y\ne0$ and $x\ge1$. If $z=0$, then
$x^3y^2=1$, giving exactly $(x,y,z)=(1,1,0),(1,-1,0)$; both satisfy the
claimed bounds because $K\ge1$.

Now put $t=|z|\ge1$ and $N=t^4+1=x^3y^2$. Each prime divisor $p$ of $N$
divides $x$ or $y$, and hence its exponent in $N$ is at least two.
Unique prime factorization therefore gives
\begin{equation}
\rad(N)^2\mid N,\qquad\text{and thus}\qquad\rad(N)\le N^{1/2}.
\label{eq:squareful}
\end{equation}
The integers $1,t^4,N$ are pairwise coprime: any common divisor of $t^4$
and $N$ divides $N-t^4=1$. Also $\rad(t^4)=\rad(t)\le t\le N^{1/4}$, so
\begin{equation}
\rad(1\cdot t^4\cdot N)
\le\rad(t^4)\rad(N)\le N^{1/4}N^{1/2}=N^{3/4}.
\label{eq:radical}
\end{equation}
Applying $abc$ with $\varepsilon=1/6$ to $1+t^4=N$ yields
\[
N\le K\,\rad(t^4N)^{7/6}\le K N^{7/8}.
\]
Since $N>0$, division gives $N^{1/8}\le K$, hence $N\le K^8$.
Finally $t^4<N$, $x^3\le N$, and $y^2\le N$ give the stated bounds.
Only finitely many integer triples lie in these bounded ranges.
\end{proof}

\newpage
\section*{2. Intuition: why the exponents force a bound}
The useful feature of $x^3y^2$ is that it has no prime factor occurring
only once. Such a positive integer is called \emph{powerful}.
Its size can grow much faster than its radical: a prime power $p^e$
contributes $p^e$ to the integer but only $p$ to the radical.
Because every exponent here is at least two, the radical uses at most
half of the size exponent, as expressed by \eqref{eq:squareful}.

The natural additive identity is
\[
\underbrace{1}_{\text{radical }1}
\; + \;
\underbrace{t^4}_{\text{radical at most }N^{1/4}}
\; = \;
\underbrace{N}_{\text{radical at most }N^{1/2}}.
\]
Thus the product of all distinct primes in this coprime triple is at most
$N^{3/4}$. The conjecture says that the sum $N$ cannot substantially exceed
that product raised to a power arbitrarily close to one. The gap
\[
\frac14+\frac12=\frac34<1
\]
is precisely what makes a finiteness proof possible.

More explicitly, any fixed $0<\varepsilon<1/3$ gives
\[
N\le K_\varepsilon N^{\frac34(1+\varepsilon)},
\qquad
N^{\frac{1-3\varepsilon}{4}}\le K_\varepsilon,
\qquad
N\le K_\varepsilon^{\frac4{1-3\varepsilon}}.
\]
The exponent on the middle left-hand side is positive. Since
$K_\varepsilon$ is independent of the solution, $N$ cannot grow without
bound. Choosing $\varepsilon=1/6$ simply makes this exponent $1/8$.

\section*{3. Scope of the conclusion}
No coprimality assumption on $x$ and $y$ is needed. The required
coprimality comes from the consecutive integers $t^4$ and $t^4+1$.
Negative $x$ and $y=0$ are impossible by positivity; both signs of $y$
and $z$ are included through absolute values.

The two solutions $(1,\pm1,0)$ are explicit. The argument does not assert
that these are the only solutions: it proves that the total solution set
is finite under $abc$. The stated form of the conjecture supplies no
numerical value for $K_{1/6}$, so the bound above is not a numerical search
cutoff and does not by itself produce a complete list.

The proof uses no external results beyond the stated conjecture and
elementary unique prime factorization. No internet search or computational
search is needed.
\end{document}
